\documentclass[11pt]{article}
\usepackage[margin=1in]{geometry}
\usepackage{multicol}
\usepackage{booktabs}

\title{Multicolumn Text Stress Case}
\author{Test Corpus}
\date{}

\begin{document}
\maketitle

\begin{abstract}
This paper-style fixture tests the \texttt{multicols} environment with ordinary
prose, section headings, and a compact list. The document remains self-contained
and avoids floating material inside the multicolumn region.
\end{abstract}

\section{Field Notes}
\label{sec:field-notes}

\begin{multicols}{2}
\subsection*{Site Context}
The study setting is a synthetic network of offices that adopted a common intake
form during the same quarter. Each office kept its local scheduling process, but
the intake fields used identical definitions. This design creates a simple
example in which textual notes need to remain readable after conversion.

\subsection*{Coding Rules}
The coding protocol asks reviewers to classify each record before inspecting
outcomes. The ordered checks are short enough for a list, but they still create
indentation and spacing decisions within the multicolumn environment:
\begin{itemize}
\item confirm that the record has a valid intake date;
\item mark whether the supervisor review field is complete;
\item flag records with inconsistent follow-up status.
\end{itemize}

\subsection*{Interpretation}
The two-column passage should preserve reading order from the left column to the
right column. Line breaks may differ after conversion, although the paragraphs
and list items should remain grouped within the multicolumn region.
\end{multicols}

\section{Check Counts}
\label{sec:check-counts}

\begin{table}[htbp]
\centering
\caption{Synthetic audit outcomes}
\label{tab:audit}
\begin{tabular}{lcc}
\toprule
Audit rule & Passed & Flagged \\
\midrule
Valid intake date & 217 & 4 \\
Complete supervisor field & 198 & 23 \\
Consistent follow-up status & 204 & 17 \\
\bottomrule
\end{tabular}
\end{table}

The table outside the multicolumn environment provides a simple reference point
for ordinary paragraph and table rendering in the same document.

\end{document}
